% !TEX root = ../main.tex
\chapter{Reductions}
\label{ch:reductions}

Reductions are operations that collapse one or more axes of a
tensor, replacing each collapsed axis with a single aggregated
value. They differ from unary and binary functions in three ways:
the output shape differs from the input shape (one or more
dimensions are removed), the operation requires a choice of which
axes to collapse, and numerical stability becomes a first-class
concern for long reductions.

\Lucid\ ships about 20 reduction ops. They all instantiate
\code{ReduceKernel<Op>}
(\chref{ch:kernel_abstraction}~Section~\ref{sec:kernel_abstraction:reduce})
with a different reducer struct, but the user-facing axis
arguments and the numerical-stability machinery are uniform.

\secref{sec:reductions:semantics} describes the axis semantics.
\secref{sec:reductions:simple} treats the simple reductions
(\code{sum}, \code{mean}, \code{max}, etc.).
\secref{sec:reductions:statistical} treats statistical reductions
(\code{var}, \code{std}, \code{logsumexp}).
\secref{sec:reductions:argument} treats argument reductions
(\code{argmax}, \code{argmin}).
\secref{sec:reductions:boolean} treats boolean reductions
(\code{any}, \code{all}).
\secref{sec:reductions:stability} treats numerical stability.

\section{Axis Semantics}
\label{sec:reductions:semantics}

Every reduction op takes an optional \code{dim} (or \code{dims})
argument and a \code{keepdims} flag:

\begin{itemize}
  \item \code{dim=None} (or omitted): reduce over all axes,
    producing a scalar.
  \item \code{dim=k} for a single integer \code{k}: reduce over
    axis \(k\), producing a tensor with one fewer dimension (or
    the same number with \code{keepdims=True}).
  \item \code{dim=(k1, k2, ...)} for a tuple: reduce over the
    listed axes, producing a tensor with that many fewer
    dimensions (or the same with \code{keepdims=True}).
\end{itemize}

Negative axis values are interpreted as offsets from the end (NumPy
convention).

\figref{fig:reductions:axis_collapse} visualises the axis-collapse
semantics for the three common cases.

\begin{figure}[t]
\centering
\begin{adjustbox}{max width=\textwidth, center}
\begin{tikzpicture}[
  cell/.style={draw, minimum width=0.6cm, minimum height=0.6cm,
               font=\footnotesize, fill=lucid.info.bg, inner sep=0pt},
  outcell/.style={draw, minimum width=0.6cm, minimum height=0.6cm,
              font=\footnotesize, fill=lucid.ok.bg, very thick,
              inner sep=0pt},
  hdr/.style={font=\small\sffamily\bfseries, align=center},
  caplbl/.style={font=\footnotesize\itshape, color=lucid.muted,
              align=center},
]
  % Input: 3x4 matrix (top-left)
  \node[hdr] at (0.9, 3) {Input shape $(3, 4)$};
  \foreach \i in {0, 1, 2} {
    \foreach \j in {0, 1, 2, 3} {
      \node[cell] at (\j*0.6, 2-\i*0.6) {};
    }
  }

  % Output 1: sum over dim=0 (row of 4 cells, top-right)
  \node[hdr] at (6, 3) {\code{sum(dim=0)} $\to$ $(4,)$};
  \foreach \j in {0, 1, 2, 3} {
    \node[outcell] at (5.1+\j*0.6, 2) {};
  }
  \node[caplbl] at (6, 1.0)
    {dim 0 collapsed:\\rank drops to 1};

  % Output 2: sum over dim=1 keepdims (column of 3 cells, bottom-left)
  \node[hdr, align=center] at (0.9, -1.0)
    {\code{sum(dim=1, keepdims=True)} $\to$ $(3, 1)$};
  \foreach \i in {0, 1, 2} {
    \node[outcell] at (0.9, -2-\i*0.6) {};
  }
  \node[caplbl] at (0.9, -4.5)
    {dim 1 collapsed but\\preserved (size-1)};

  % Output 3: global sum -> scalar
  \node[hdr] at (6, -1.0) {\code{sum()} $\to$ scalar};
  \node[outcell] at (6, -2) {};
  \node[caplbl] at (6, -2.9)
    {all dims collapsed};
\end{tikzpicture}
\end{adjustbox}
\caption[Reduction shape transformations.]{Reduction along
different axis specifications applied to the same $(3, 4)$ input.
Without \code{keepdims}, the reduced dimensions disappear (top
right). With \code{keepdims=True}, they remain as size-1
placeholders (bottom left), which makes the result naturally
broadcast-compatible with the original input --- the pattern used
throughout normalisation and centring code. A global reduction
(bottom right) collapses every dimension and returns a scalar.}
\label{fig:reductions:axis_collapse}
\end{figure}

\subsection{The keepdims flag}
\label{ssec:reductions:keepdims}

\code{keepdims=True} preserves the rank by leaving a 1 in the
reduced dimension's slot. For \code{x} of shape \((B, C, H, W)\),
\code{x.sum(dim=1, keepdims=True)} returns shape \((B, 1, H, W)\)
whereas \code{x.sum(dim=1)} returns \((B, H, W)\).

The keepdims form is the natural input to a subsequent broadcast:
\code{x - x.mean(dim=-1, keepdims=True)} subtracts the per-row
mean without needing an explicit \code{unsqueeze}. This pattern
recurs throughout neural-network code (layer normalisation, mean
centring).

\subsection{The \texttt{dim} vs \texttt{axis} naming}
\label{ssec:reductions:dim_axis}

\Lucid\ standardised on \code{dim} (singular) and \code{dims}
(plural) in the engine-wide rename that happened during Phase 7
(vault: \code{retro-axis-to-dim-rename}). The change brought parity
with the \refframework\ convention and away from NumPy's \code{axis}
naming. The Python dispatchers translate \code{axis=...} to
\code{dim=...} for NumPy-style callers, so both spellings work; the
canonical form is \code{dim}.

\section{Simple Reductions}
\label{sec:reductions:simple}

The simple reductions:

\begin{itemize}
  \item \code{sum(x)}: \(\sum_i x_i\). Backward: gradient
    broadcasts to every input position.
  \item \code{prod(x)}: \(\prod_i x_i\). Backward: gradient times
    product divided by each entry (with the standard zero-input
    workaround).
  \item \code{mean(x)}: \(\frac{1}{N}\sum_i x_i\). Backward:
    gradient divided by \(N\), broadcast.
  \item \code{max(x)}: \(\max_i x_i\). Backward: gradient routes
    to the maximum-valued input; ties get a 1/k split where \(k\)
    is the count of maxima.
  \item \code{min(x)}: analogous.
  \item \code{amax(x), amin(x)}: aliases for \code{max} and
    \code{min} (NumPy compatibility).
  \item \code{median(x)}: median value; backward: gradient routes
    to the median-rank element with the 1/k tie-breaking.
  \item \code{quantile(x, q)}: linear-interpolation quantile.
  \item \code{nanmean, nansum, nanmedian}: like the
    above but skip NaN entries.
\end{itemize}

The variants prefixed \code{nan} skip NaN inputs in the reduction.
The implementation is a mask-then-reduce: replace NaNs with the
reducer's identity, reduce, then renormalise where appropriate.

\subsection{prod's backward at zero}
\label{ssec:reductions:prod_zero}

\code{prod}'s backward is mathematically \(\partial_i \prod_j x_j =
\prod_{j \neq i} x_j\). The naive computation as
\code{prod(x) / x[i]} is undefined when any \code{x[i]} is zero.
\Lucid\ uses the standard workaround: explicitly compute
\(\prod_{j \neq i}\) without the division by tracking partial
products from both sides of \(i\). The implementation is in
\code{ops/gfunc/Reduction.cpp}.

\section{Statistical Reductions}
\label{sec:reductions:statistical}

The statistical family:

\begin{itemize}
  \item \code{var(x, unbiased=True)}: variance, optionally with
    Bessel's correction (\(N-1\) denominator).
  \item \code{std(x, unbiased=True)}: standard deviation.
  \item \code{nanvar, nanstd}: NaN-skipping variants.
  \item \code{logsumexp(x)}: \(\log\sum_i e^{x_i}\), computed
    via the max-shift trick.
  \item \code{cumsum(x), cumprod(x)}: cumulative reductions
    (technically a scan, not a reduction; output has the same
    shape as input).
  \item \code{cummax(x), cummin(x)}: cumulative max/min.
\end{itemize}

\subsection{Variance and Bessel's correction}
\label{ssec:reductions:bessel}

The \code{unbiased=True} flag divides by \(N - 1\) instead of \(N\),
producing an unbiased estimator of the population variance. The
default is \code{True}, matching \refframework. NumPy's default
\code{ddof=0} (biased) is the opposite, which has surprised users
porting code; we document this prominently in the API reference.

The historical schema migration \code{batch\_norm v1$\to$v2}
(\chref{ch:op_registry}) was precisely this issue: BN's running
variance was stored biased and changed to unbiased in v2 to match
\refframework. The migration rescales by \(N/(N-1)\).

\subsection{logsumexp and the max-shift trick}
\label{ssec:reductions:logsumexp}

Naive \code{log(sum(exp(x)))} overflows for moderately large
\(x\) (e.g., \code{exp(1000)} is \(+\infty\) in any IEEE 754
precision). The standard workaround:

\[
\text{logsumexp}(x) = m + \log \sum_i e^{x_i - m},
\quad m = \max_i x_i.
\]

The subtraction inside \code{exp} keeps the exponents bounded, and
the added \(m\) recovers the correct magnitude. \Lucid's
\code{logsumexp} kernel implements this directly rather than
chaining \code{log}-\code{sum}-\code{exp}.

The backward of \code{logsumexp} is the softmax:
\(\partial_i \text{logsumexp}(x) = \text{softmax}(x)_i\). This is
the gradient identity that makes softmax-cross-entropy losses
numerically tractable.

\section{Argument Reductions}
\label{sec:reductions:argument}

\begin{itemize}
  \item \code{argmax(x, dim)}: index of the maximum along
    \code{dim}.
  \item \code{argmin(x, dim)}: index of the minimum.
  \item \code{topk(x, k, dim, largest=True)}: returns the top
    \code{k} values and their indices along \code{dim}.
  \item \code{kthvalue(x, k, dim)}: returns the \(k\)-th value (in
    sorted order) along \code{dim}.
\end{itemize}

Argument reductions return integer tensors and are not
differentiable. Backward returns zero gradient.

\subsection{The topk backward through value}
\label{ssec:reductions:topk_backward}

The \emph{values} returned by \code{topk} (not the indices)
\emph{are} differentiable: the gradient routes to the corresponding
input positions identified by the indices. \Lucid's \code{topk}
implementation handles this in two ways depending on use:

\begin{itemize}
  \item If the caller uses only the values
    (\code{vals, \_ = lucid.topk(x, k)} and then \code{vals.sum().backward()}),
    the backward scatters the gradient to the top-\(k\) input
    positions.
  \item If the caller uses only the indices, no gradient is
    constructed.
\end{itemize}

The autograd engine determines this by inspecting which output is
referenced downstream.

\section{Boolean Reductions}
\label{sec:reductions:boolean}

\begin{itemize}
  \item \code{any(x)}: true if any entry is true (or nonzero).
  \item \code{all(x)}: true if every entry is true.
\end{itemize}

Boolean reductions are not differentiable. They are commonly used
in conditional logic that does not need gradients (early-stopping
checks, NaN detection in the gradients themselves).

\section{Numerical Stability}
\label{sec:reductions:stability}

Long reductions are where floating-point error compounds the most,
and \Lucid\ has several mechanisms for managing it.

\subsection{Reduction-in-higher-precision}
\label{ssec:reductions:higher_precision}

Reductions over FP16 or BF16 tensors automatically accumulate in
FP32, then cast back. The kernel template handles this through the
\code{AmpPolicy::ForceFloat32} declaration on the reduction op's
schema (\chref{ch:op_registry}).

Without this, summing \(10^6\) FP16 values would lose roughly
\(\log_2(10^6) \approx 20\) bits of precision in the accumulator,
which is much more than FP16's \(\sim 11\)-bit mantissa. With FP32
accumulation, the loss is \(\sim 20\) bits in FP32's 23-bit
mantissa, which leaves \(\sim 3\) bits of precision --- still
enough to be useful, and the result is cast back to FP16 only at
the end.

\subsection{Kahan summation}
\label{ssec:reductions:kahan}

\Lucid\ does not implement Kahan-style compensated summation. The
trade-off: Kahan summation gives much better precision but at
roughly \(4\times\) the cost. For our typical workloads, the
FP32-accumulation approach gives enough precision without the cost.

The decision is documented in the FP32-accumulation schema
declaration; a user who needs better precision can request it by
casting to FP64 (CPU only) before the reduction.

\subsection{Compensated mean: sum then divide}
\label{ssec:reductions:mean_div}

\code{mean(x)} is computed as \code{sum(x) / numel(x)} rather than
as an incremental running mean. The reasoning: \code{sum} is
implemented in higher precision (FP32) and the division by a
known integer is a single round-to-nearest operation. The
incremental running mean accumulates roundoff at every step,
whereas the divide-once approach accumulates it only at the
final step.

This is the optimisation referenced in the no-compile sweep
(\chref{ch:no_compile_sweep}): replacing \code{mean} with
\code{sum / N} folds the divide earlier in the computation graph,
which can yield a measurable performance improvement on
bandwidth-bound workloads.

\section{Performance Characteristics}
\label{sec:reductions:perf}

\subsection{Contiguous vs strided axis}
\label{ssec:reductions:contig_strided}

Reducing a contiguous axis is much faster than reducing a strided
one: the contiguous case allows the backend to issue a single
strided-load instruction per cache line, while the strided case
requires per-element address computation. The kernel template
chooses between two implementations based on the input's strides
(\chref{ch:kernel_abstraction}~Section~\ref{ssec:kernel_abstraction:axis_select}).

\subsection{The MPSGraph carve-out: LayerNorm}
\label{ssec:reductions:layernorm_mps}

Although individual reductions do not have an \MPSGraph\ carve-out
(\MLX's reductions are competitive), \code{layer\_norm} ---  which
is composed of \code{mean} and \code{var} reductions followed by a
broadcast affine --- does. The \MPSGraph\ implementation fuses
all three steps into a single kernel and achieves substantially
better throughput than the equivalent \MLX\ sequence.

The carve-out is per-op (it triggers for the composite
\code{layer\_norm}, not for the individual reductions). See
\chref{ch:mpsgraph_dispatch} for the full list.

\section{Worked Numerical-Stability Examples}
\label{sec:reductions:worked_stability}

The general principle of FP32-accumulation under AMP is best
understood through worked examples.

\subsection{Example: large-batch mean in FP16}
\label{ssec:reductions:large_mean}

Consider computing the mean of $10^6$ FP16 values, each drawn from
a normal distribution with variance 1. Naively in FP16:

\begin{lstlisting}[style=lucid,language=LucidPython]
import lucid
x = lucid.randn(1_000_000, dtype=lucid.float16)
m = x.mean()                   # FP16 accumulator: WRONG
\end{lstlisting}

FP16 has 11-bit mantissa, giving roughly 4-decimal-digit precision.
After accumulating $10^6$ random values of magnitude $\sim 1$, the
running sum reaches $\sim 1000$ (because $\sqrt{10^6} = 1000$ by
central limit). Each addition of a new $\sim 1$-magnitude value
into a $\sim 1000$ accumulator loses the lowest bits; the
accumulated error is on the order of $10^6 \cdot 10^{-4} \cdot 1 =
100$. The reported mean is wrong by a factor of $\sim 100$.

\Lucid's actual implementation forces FP32 accumulation:

\begin{lstlisting}[style=lucid,language=LucidPython]
x = lucid.randn(1_000_000, dtype=lucid.float16)
m = x.mean()                   # FP32 accumulator (per AmpPolicy)
                               # then cast back to FP16
\end{lstlisting}

FP32's 23-bit mantissa gives roughly 7-decimal-digit precision.
Accumulating $10^6$ values of magnitude $1$ into a $1000$-magnitude
accumulator loses $\sim 10^{-7}$ per add, total error $\sim 0.1$,
which is well within FP16's representable precision for a result
of magnitude $\sim 0$. The cast back to FP16 at the end loses no
additional precision because the result already fits.

\subsection{Example: softmax with extreme logits}
\label{ssec:reductions:softmax_extreme}

Softmax is defined as $\text{softmax}(x)_i = e^{x_i} / \sum_j
e^{x_j}$. With extreme inputs, the naive computation overflows:
$e^{1000}$ is $+\infty$ in FP32.

The standard trick subtracts the max: $\text{softmax}(x)_i =
e^{x_i - m} / \sum_j e^{x_j - m}$ where $m = \max_i x_i$. The
exponents are now all in $(-\infty, 0]$ and the largest is
exactly 0, so the largest $e^{x_i - m} = 1$ and the sum is
between 1 and $N$. No overflow possible.

\Lucid's \code{softmax} primitive implements this internally;
\code{lucid.nn.functional.softmax} delegates. The composite
implementation (\chref{ch:composite_ops}) makes the trick
explicit.

\subsection{Example: variance with Bessel correction}
\label{ssec:reductions:var_bessel}

For a sample $x_1, \ldots, x_N$, the unbiased variance estimator
is
\[
s^2 = \frac{1}{N - 1} \sum_{i=1}^N (x_i - \bar{x})^2.
\]
\Lucid's \code{var} defaults to \code{unbiased=True}, matching
\refframework. NumPy's default (\code{ddof=0}) uses $N$ instead of
$N-1$, the biased estimator.

The two differ by a factor of $N/(N-1)$, which is small for large
$N$ but matters for small samples. The
\code{batch\_norm v1$\to$v2} migration was exactly this issue:
running variance was stored biased and changed to unbiased to
match \refframework's BN; checkpoints saved before the migration
need rescaling.

\section{Performance and Benchmarks}
\label{sec:reductions:perf_benchmarks}

\subsection{Bandwidth ceiling}
\label{ssec:reductions:bw_ceiling}

For large tensors a reduction is bandwidth-bound: read $N$
elements, write 1 or a few output elements. The throughput ceiling
on M3~Max (400~GB/s) for FP32 reductions is approximately
$400 / 4 = 100$ giga-elements per second. Measured rates for the
contiguous case come within 70--85\,\% of this ceiling depending
on the reduction op and the axis structure.

For small tensors the launch overhead dominates: a reduction over
a 1~KB tensor takes the same wall-clock time as one over a 64~KB
tensor on the GPU, because both pay the kernel-launch cost.

\subsection{Strided-axis penalty}
\label{ssec:reductions:strided_penalty}

Reducing along an inner axis with non-unit stride is roughly
$2$--$4\times$ slower than reducing along a unit-stride axis. The
penalty comes from address computation and from poor cache
utilisation (the strided reads do not coalesce into cache-line
loads).

\tabref{tab:reductions:strided_bench} reports measured throughput
on M3~Max for a $1024 \times 1024 \times 16$ FP32 tensor under
different reduction axes.

\begin{table}[t]
\centering
\small
\begin{tabular}{llrl}
\toprule
Reduction & Axis & Throughput (GE/s) & Cache behaviour \\
\midrule
\code{sum} & dim=2 (innermost, contig) & 82 & coalesced loads \\
\code{sum} & dim=0 (outermost, contig per row) & 75 & coalesced loads \\
\code{sum} & dim=1 (middle, contig per row) & 68 & coalesced loads \\
\code{sum} & dim=1 (middle, after transpose) & 22 & strided loads \\
\code{sum} & dim=-1 (innermost, after transpose) & 18 & strided + scatter \\
\bottomrule
\end{tabular}
\caption[Reduction throughput by axis.]{\code{sum} throughput on
M3~Max for a $1024 \times 1024 \times 16$ FP32 tensor under
different axis configurations. The strided cases (last two rows)
are 3--4$\times$ slower than the contiguous cases. The remedy is
to either \code{contiguous()} the tensor before the reduction (one
extra pass over the data) or to fuse the reduction with the
upstream transpose at the kernel level, which \Lucid's dispatcher
does opportunistically when the pattern is recognisable.}
\label{tab:reductions:strided_bench}
\end{table}

\section{Summary and Forward References}
\label{sec:reductions:summary}

The reduction category contains roughly 20 ops covering simple
reductions (sum, mean, max, etc.), statistical reductions (var,
std, logsumexp), argument reductions (argmax, topk, kthvalue), and
boolean reductions (any, all). All instantiate
\code{ReduceKernel<Op>} and inherit the kernel's axis-iteration
machinery. Numerical stability is managed through FP32 accumulation
under AMP and through algorithm-specific tricks for the
exponential-family reductions (logsumexp's max shift).

The next chapter, \chref{ch:shape_view}, treats shape and view
manipulation --- operations that change a tensor's shape without
changing its element values. \chref{ch:indexing} treats indexing
and gather/scatter. \chref{ch:composite_ops} closes Part III with
the pure-Python composite layer.

\begin{lucidnote}
For the user: reductions are exposed both as free functions
(\code{lucid.sum(x)}) and as Tensor methods (\code{x.sum()}). The
common patterns --- per-row mean with \code{dim=-1, keepdims=True},
batch-wise reduction with \code{dim=0} --- are idiomatic and run
without any explicit per-element loop.
\end{lucidnote}
